% ============================================================
%  本文档为「人机交互导论」课程设计报告，沿用 TongjiThesis 模板格式
% ============================================================

% ============================================================
%  封面信息
% ============================================================
\school{计算机科学与技术学院}
\major{计算机科学与技术}
% 小组作业：4 名成员，姓名与学号用顿号分隔，左右一一对应
\student{学号一、学号二、学号三、学号四}{姓名一、姓名二、姓名三、姓名四}
\thesistitle{Mood Radio：面向情绪感知的本地音乐电台}{多模态人机交互系统的设计与实现}
\thesistitleeng{Mood Radio: An Emotion-Aware Local Music Radio}{Design and Implementation of a Multimodal HCI System}
\thesisadvisor{倪张凯}
\thesisdate{2026}{6}{14}

% ============================================================
%  信息说明页
% ============================================================
\infotype{thesis}

\infoabstract{%
本课程设计面向「音乐推荐难以理解用户当下情绪」这一痛点，设计并实现了一个本地优先的情绪感知音乐电台 \mbox{Mood Radio}。系统以 Valence--Arousal 二维情绪模型为公共语义，将摄像头表情、手动选择、自然语言聊天、时段推断与头部手势五种通道统一映射到目标情绪坐标，并基于本地曲库进行复合打分推荐（情绪距离、用户反馈、新鲜度与时间惩罚加权融合）。前端采用 React + Vite + TailwindCSS 构建，后端使用 \mbox{C++17} + cpp-httplib 提供 25 个以上的 REST 接口，配合 SQLite 与 Python 离线分析脚本完成音乐情绪建模。系统强调隐私优先（摄像头与音乐均本地处理、LLM 密钥不落盘）、可解释推荐（情绪图谱与评分徽标）以及沉浸式视觉反馈（Aurora 流动背景、情绪驱动主题、黑胶旋转封面）。课堂模拟评估显示，9/10 用户认可本地隐私设计，8/10 用户能理解推荐逻辑，本项目较好地体现了多模态、情境感知、可解释与隐私保护等关键 HCI 设计原则。
}

\infothesiswords{11000}

\infomaterials{
  \item 项目源代码仓库（hyx\_feature 分支）；
  \item 启动脚本 start.ps1 与依赖说明。
}
